\documentclass[11pt,a4paper]{article}
\usepackage[utf8]{inputenc}
\usepackage[french]{babel}
\usepackage[T1]{fontenc}
\usepackage{geometry}
\geometry{top=2.5cm, bottom=2.5cm, left=2.5cm, right=2.5cm}
\usepackage{amsmath,amssymb}
\usepackage{graphicx}
\usepackage{xcolor}
\usepackage{hyperref}
\usepackage{listings}
\usepackage{booktabs}
\usepackage{fancyhdr}
\usepackage{tikz}
\usetikzlibrary{shapes.geometric, arrows, positioning}

\definecolor{primaryblue}{RGB}{24, 76, 120}
\definecolor{secondaryblue}{RGB}{41, 128, 185}
\definecolor{darkgray}{RGB}{50, 50, 50}
\definecolor{codebg}{RGB}{245, 247, 250}
\definecolor{codeframe}{RGB}{210, 215, 225}

\hypersetup{
    colorlinks=true,
    linkcolor=primaryblue,
    urlcolor=secondaryblue,
    titlecolor=primaryblue
}

\pagestyle{fancy}
\fancyhf{}
\rhead{\textcolor{darkgray}{\small Rapport de Projet -- Warehouse Fusion}}
\lhead{\textcolor{darkgray}{\small Robotique \& Apprentissage par Renforcement}}
\rfoot{\thepage}

\lstset{
    backgroundcolor=\color{codebg},
    frame=single,
    rulecolor=\color{codeframe},
    basicstyle=\ttfamily\small,
    keywordstyle=\color{primaryblue}\bfseries,
    commentstyle=\color{gray}\itshape,
    stringstyle=\color{teal},
    breaklines=true,
    captionpos=b,
    showstringspaces=false
}

\begin{document}

% ── PAGE DE GARDE ──────────────────────────────────────────────────────────
\begin{titlepage}
    \centering
    \vspace*{1cm}
    
    {\Huge \bfseries \color{primaryblue} WAREHOUSE FUSION}\\[0.4cm]
    {\Large \bfseries Système Unifié de Navigation Autonome et de Manipulateur Robotique par Apprentissage par Renforcement}\\[1.5cm]
    
    \rule{\linewidth}{1.5pt}\\[0.8cm]
    {\Large \textbf{Rapport Technique Intégral}}\\[0.3cm]
    {\large Intégration de ROS 2 Nav2, Robosuite, MuJoCo, PyTorch et Stable-Baselines3}\\[0.8cm]
    \rule{\linewidth}{1.5pt}\\[2.5cm]
    
    \begin{minipage}{0.48\textwidth}
        \begin{flushleft} \large
            \textbf{Auteur :} Rabeb\\
            \textbf{Projet :} Warehouse Fusion\\
            \textbf{Robot :} Franka Emika Panda \& Mobile Trailer
        \end{flushleft}
    \end{minipage}
    \begin{minipage}{0.48\textwidth}
        \begin{flushright} \large
            \textbf{Environnement :} Linux (Ubuntu 22.04)\\
            \textbf{Framework :} ROS 2 Humble / Robosuite\\
            \textbf{Date :} Septembre 2026
        \end{flushright}
    \end{minipage}
    
    \vfill
    {\small \color{gray} Rapport de développement, d'entraînement RL et d'intégration système}
\end{titlepage}

\newpage
\tableofcontents
\newpage

% ── SECTION 1 : RESUME EXECUTIF & OBJECTIFS ────────────────────────────────
\section{Résumé Exécutif et Objectifs du Projet}

Le projet \textbf{Warehouse Fusion} s'inscrit dans le domaine de la robotique d'entrepôt de nouvelle génération. L'objectif principal est d'orchestrer la coopération autonome entre un **robot mobile à remorque (Trailer)** gérant la logistique horizontale et un **bras manipulateur (Franka Emika Panda)** assurant la saisie et le transbordement d'objets (Pick \& Place).

\subsection{Objectifs Principaux}
\begin{enumerate}
    \item \textbf{Navigation Autonome (Nav2)} : Guider le chariot mobile à travers l'entrepôt vers la station de chargement du bras Panda en évitant les obstacles et en assurant un amarrage précis ($x = -0.24\text{m}, y = 2.80\text{m}$).
    \item \textbf{Apprentissage par Renforcement (RL)} : Entraîner une politique sous \textbf{Robosuite} à l'aide de l'algorithme \textbf{PPO (Proximal Policy Optimization)} pour permettre au bras Panda d'intercepter, d'attraper et de soulever l'objet cible (une canette rouge).
    \item \textbf{Orchestration Unifiée} : Connecter le sous-système de navigation et le sous-système de manipulation via des signaux de synchronisation ROS 2 (`/trailer/docking_status` et `/panda/task_completed`).
    \item \textbf{Validation Expérimentale} : Valider l'efficacité de l'agent RL en mesurant l'évolution de la courbe de récompense (passant de $0.89$ à $14.34$) et assurer une solution unifiée 100\% fonctionnelle.
\end{enumerate}

% ── SECTION 2 : ARCHITECTURE SYSTEME & STACKS ──────────────────────────────
\section{Architecture Système et Stacks Technologiques}

L'architecture repose sur deux piliers technologiques complémentaires : l'écosystème ROS 2 pour la communication distribuée et le framework Robosuite/MuJoCo pour la simulation physique et l'apprentissage RL.

\begin{table}[h!]
\centering
\caption{Synthèse des Stacks et Outils Logiciels Utilisés}
\begin{tabular}{lll}
\toprule
\textbf{Composant} & \textbf{Technologie / Bibliothèque} & \textbf{Rôle dans le projet} \\
\midrule
Système d'Exploitation & Linux Ubuntu 22.04 LTS & OS hôte pour la simulation et ROS 2 \\
Middleware Robotique & ROS 2 Humble Hawksbill & Orchestration, TF2, Topics \& Actions \\
Navigation Autonome & Nav2 (Planner, Controller, BT) & Planification de trajectoire du chariot \\
Localisation \& TF & `robot_localization` (EKF) & Fusion de données IMU/GPS/Odométrie \\
Moteur de Simulation & MuJoCo / Robosuite 1.5.2 & Physique de contact et dynamique du bras \\
Apprentissage RL & Stable-Baselines3 & Implémentation des algorithmes PPO et SAC \\
Calcul Numérique & PyTorch (Version CPU) & Réseaux de neurones (Acteur-Critique) \\
Visualisation & RViz2 \& GLFW Renderer & Visualisation 3D et suivi télémétrique \\
\bottomrule
\end{tabular}
\end{table}

\subsection{Chaîne de Transformations TF2}
La synchronisation spatiale est garantie par la chaîne TF2 complète :
$$\text{world} \longrightarrow \text{map} \longrightarrow \text{odom} \longrightarrow \text{base\_footprint} \longrightarrow \text{trailer\_cart\_drop\_point}$$

Le repère racine du bras manipulateur (`world`) est lié avec le repère de navigation Nav2 (`map`) via un transformateur statique rigide à l'identité.

% ── SECTION 3 : INSTALLATION ET CONFIGURATION ──────────────────────────────
\section{Pipeline d'Installation et Configuration du Workspace}

Pour garantir l'isolation et la reproductibilité des environnements, un environnement virtuel Python dédié a été configuré dans `~/Downloads/robosuite-master/venv`.

\subsection{Commandes d'Installation et de Dépendances}
\begin{lstlisting}[language=bash, caption={Installation des packages et du venv}]
# 1. Activation du Venv
source /home/rabeb/Downloads/robosuite-master/venv/bin/activate

# 2. Installation de PyTorch CPU et Stable-Baselines3
pip install torch stable-baselines3 tensorboard \
    --index-url https://download.pytorch.org/whl/cpu \
    --extra-index-url https://pypi.org/simple

# 3. Configuration des Macros Robosuite
python /home/rabeb/Downloads/robosuite-master/robosuite/scripts/setup_macros.py

# 4. Compilation du workspace ROS 2
cd /home/rabeb/warehouse_fusion_ws
colcon build --packages-select warehouse_orchestrator
source install/setup.bash
\end{lstlisting}

% ── SECTION 4 : ALGORITHMES ET METHODES ─────────────────────────────────────
\section{Algorithmes et Stratégies d'Apprentissage (RL vs Déterministe)}

Deux stratégies ont été développées et évaluées pour la tâche de \textit{Pick \& Place}.

\subsection{1. Algorithme Proximal Policy Optimization (PPO)}
L'agent RL apprend une politique stochastique $\pi_\theta(a|s)$ représentée par un réseau de neurones multilayer perceptron (MLP). L'objectif optimisé par PPO est donné par :
$$L^{CLIP}(\theta) = \hat{\mathbb{E}}_t \left[ \min\left(r_t(\theta)\hat{A}_t, \, \text{clip}(r_t(\theta), 1-\epsilon, 1+\epsilon)\hat{A}_t\right) \right]$$
où $r_t(\theta) = \frac{\pi_\theta(a_t|s_t)}{\pi_{\theta_{old}}(a_t|s_t)}$ est le ratio de probabilité et $\hat{A}_t$ est l'avantage estimé par le réseau Critique.

\subsubsection{Résultats Expérimentaux de l'Entraînement}
Le modèle a été entraîné sur **200 000 timesteps** (durée : 32.66 minutes sur CPU).

\begin{table}[h!]
\centering
\caption{Évolution des Métriques RL entre 50k et 200k pas d'apprentissage}
\begin{tabular}{lccc}
\toprule
\textbf{Métrique} & \textbf{50 000 Pas} & \textbf{200 000 Pas} & \textbf{Interprétation} \\
\midrule
Récompense moyenne (`ep_rew_mean`) & $0.89$ & \textbf{$5.28$} & Hausse spectaculaire des performances \\
Récompense maximale observée & $1.86$ & \textbf{$14.34$} & Saisie et levage réussis \\
Variance expliquée (`explained_variance`) & $-0.289$ & \textbf{$+0.711$} & Excellente prédiction de la fonction de valeur \\
Écart-type des actions (`std`) & $0.917$ & \textbf{$0.630$} & Mouvements ciblés et réduction du bruit \\
Fréquence d'entraînement (`fps`) & $104$ & \textbf{$102$} & Stabilité du moteur MuJoCo \\
\bottomrule
\end{tabular}
\end{table}

\subsection{2. Contrôleur Déterministe à Machine à Etats (IK)}
En complément du modèle RL, un contrôleur cinématique déterministe à 8 états a été conçu pour garantir un taux de succès de 100\% lors de la démonstration industrielle :
\begin{enumerate}
    \item \texttt{MOVE\_ABOVE} : Alignement $XY$ au-dessus de la canette rouge ($z + 0.15\text{m}$).
    \item \texttt{DESCEND} : Descente verticale à la hauteur de saisie ($z - 0.01\text{m}$).
    \item \texttt{GRASP} : Maintien de la position et fermeture de la pince ($a_{gripper} = +1.0$).
    \item \texttt{LIFT} : Soulevé vertical de sécurité ($z + 0.22\text{m}$).
    \item \texttt{MOVE\_TO\_BIN} : Translation horizontale vers le bac du chariot amarré.
    \item \texttt{LOWER} : Descente contrôlée dans le bac ($z_{bin} + 0.10\text{m}$).
    \item \texttt{RELEASE} : Ouverture de la pince ($a_{gripper} = -1.0$) et libération.
    \item \texttt{DONE} : Tâche achevée, envoi du signal `/panda/task_completed`.
\end{enumerate}

% ── SECTION 5 : INTEGRATION GLOBALE ET RESULTATS ───────────────────────────
\section{Intégration Globale et Validation Expérimentale}

L'intégration unifiée a été concrétisée par le script d'orchestration globale :
$$\texttt{run\_full\_warehouse\_simulation.py}$$

\subsection{Validation de la Mission Complète}
Lors de l'exécution finale, la séquence d'événements s'est déroulée sans erreur :
\begin{itemize}
    \item \textbf{PHASE 1 (Nav2)} : Le chariot autonome quitte le départ entrepôt ($0.0, 0.0$), franchit le corridor principal ($-0.5, 1.2$), puis atteint la pose d'amarrage ($-0.24\text{m}, 2.80\text{m}$).
    \item \textbf{PHASE 2 (Signalisation)} : Émission de la notification `/trailer/docking_status = True`.
    \item \textbf{PHASE 3 (Pick \& Place)} : Le robot Panda détecte la canette rouge aux coordonnées $[0.076, -0.098, 0.860]$, exécute la saisie et transborde l'objet vers le bac du chariot $[0.10, 0.28, 0.80]$.
    \item \textbf{PHASE 4 (Succès)} : Émission du signal `/panda/task_completed = True` et affichage du message :
    $$\textbf{ MISSION RÉUSSIE À 100\% ! L'objet rouge est déposé dans le chariot.}$$
\end{itemize}

% ── SECTION 6 : CONCLUSION ET PERSPECTIVES ──────────────────────────────────
\section{Conclusion et Perspectives}

Le projet \textbf{Warehouse Fusion} a démontré avec succès l'intégration synergique de la navigation autonome sous ROS 2 et de la manipulation robotique avancée par Apprentissage par Renforcement et logique déterministe sous Robosuite. 

Les résultats obtenus (récompense RL grimpant à $14.34$, `explained_variance` de $0.711$, et simulation unifiée 100\% fonctionnelle) constituent une base solide pour des développements futurs :
\begin{itemize}
    \item \textbf{Sim-to-Real Transfer} : Déploiement des politiques RL entraînées sur un véritable bras Franka Emika Panda physique.
    \item \textbf{Vision Basée sur la Caméra (RGB-D)} : Entraînement de l'agent directement sur des observations visuelles (`use_camera_obs=True`).
    \item \textbf{Algorithme SAC (Soft Actor-Critic)} : Exploration de l'efficacité spectrale de SAC sur des horizons d'entraînement plus longs ($500\,000+$ pas).
\end{itemize}

\end{document}
